\scanpage{15}

Je te signale que j'ai réfléchi aux Gr-champs sur $X$. Si $G$ est un Groupe sur $X$; $N$ un $G$-Module les Gr-champs sur $X$ ``épinglés par $G,N$'' forment a priori une $2$-catégori et même une $2$-catégorie de Picard stricte, grâce à l'opération évident à la Baer. On trouve que le complexe \inserted{(de cochaines)} tronqué à l'échelon $2$ qui leur correspond est le tronqué
\[
\tau_{\leq 2}\bigl(\mathbb R\Gamma(B_{G}\mathbin{\mathrm{mod}}X,N)[1]\bigr)
\]
(\uline{NB} la cohomologie de $\mathbb R\Gamma(B_{G}\mathbin{\mathrm{mod}}X,N)$ commence en degré $1$) Plus géométr\unsure{i}quement, un Gr-champ sur $X$ éping \inserted{par $(G,N)$} est essentiellement ``la même chose'' qu'une $2$-gerbe sur $B_{G}$, liée par $N$, et \deleted{\unsure{épi}}munie d'une trivialisati au dessus de $X\approx B_{\unsure{e|c}}=(B_{G})_{/P}$ (où $P$ est le \deleted{\unsure{``|s|$\ast$}}objet de $B_{G}$ «torseur universel sous $G$»). Ces $2$-gerbes forment en fait une $3$-catégorie de Picard a priori, mais il se
% archive sheet number: 14
